\documentclass[UTF8,a4paper,11pt]{ctexart}

\usepackage{amsmath,amssymb,bm}
\usepackage{geometry}
\usepackage{hyperref}
\geometry{margin=2.2cm}
\hypersetup{
  colorlinks=true,
  linkcolor=blue,
  urlcolor=blue
}

\title{高精度数据处理中的方法理论简介\\（外推、误差传递、拟合与统计平均）}
\author{}
\date{}

\begin{document}
\maketitle

\section{符号约定与基本假设}

本文仅给出方法层面的理论简介，不涉及任何具体软件界面或实现细节。

\subsection{数据与不确定度}

设观测量为 $x$，其标准不确定度（standard uncertainty）为 $\sigma_x$。在误差传递与加权拟合中，默认假设不同输入量之间相互独立（不相关）；若存在相关性，应使用协方差矩阵形式（见后文）。

\subsection{序列与极限}

外推（extrapolation）处理的典型对象是一个收敛序列 $\{s_n\}$（例如随基组/截断等级增大的能量序列），目标是估计极限 $s_\infty=\lim_{n\to\infty}s_n$，并给出一个可信的不确定度估计。

\section{序列外推方法（Extrapolation）}

\subsection{三点外推（默认三点公式）}

给定三项 $A,B,C$（通常理解为逐步改进的三次近似），一种常用的三点外推形式为
\begin{equation}
V = \frac{(C-B)^2}{B-A} + C,
\end{equation}
其中 $V$ 为外推值。该形式的特点：
\begin{itemize}
  \item 仅需三点，无需额外模型参数；
  \item 对数据的尺度平移不敏感，但当 $B\approx A$ 时会产生数值不稳定（分母接近零）；
  \item 适合用作“快速、保守”的三点极限估计，但不保证在所有类型序列上严格成立。
\end{itemize}

不确定度估计可采用保守的“与输入点差异上界”：
\begin{equation}
U=\max\left(|V-A|,|V-B|,|V-C|\right).
\end{equation}
该 $U$ 更像是“外推与已有数据的一致性差异”，而非统计意义上的标准差。

\subsection{幂律模型外推（Power-law extrapolation）}

对许多数值近似（尤其是具有“截断阶数/基组大小”含义的序列），常见的误差模型为
\begin{equation}
E(x) = E_\infty + a\,x^{-p},
\end{equation}
其中 $x$ 表示控制精度的参数（例如基组等级或截断指标），$E_\infty$ 为目标极限，$a,p$ 为待定常数。若给定三组 $(x_1,E_1),(x_2,E_2),(x_3,E_3)$，可通过消元求解 $p$ 并得到 $E_\infty$：
\begin{align}
R &= \frac{E_1-E_2}{E_2-E_3},\\
\frac{x_1^{-p}-x_2^{-p}}{x_2^{-p}-x_3^{-p}} &= R,\\
a &= \frac{E_1-E_2}{x_1^{-p}-x_2^{-p}},\qquad
E_\infty = E_1 - a\,x_1^{-p}.
\end{align}
特点与注意事项：
\begin{itemize}
  \item 模型具有明确物理/数值动机（主导误差项为幂律衰减）；
  \item 对三点数据较敏感：若 $E_2\approx E_3$ 或 $x_2^{-p}\approx x_3^{-p}$，会导致求解不稳定；
  \item 可在 $p$ 已知/可约束时提高鲁棒性；也可通过多点拟合扩展为过定约束问题。
\end{itemize}

一种常见的保守不确定度估计是与参考点（通常取最高等级点 $E_3$ 或用户指定列）做差：
\begin{equation}
U = |E_\infty - E_{\mathrm{ref}}|.
\end{equation}

\subsection{Richardson 外推 / 序列加速}

Richardson 外推的核心思想：若近似 $S(h)$ 具有渐近展开
\begin{equation}
S(h)=S + c\,h^p + \mathcal{O}(h^{p+1}),
\end{equation}
则可组合不同步长 $h_1,h_2$ 的近似以消去主导误差项，从而提高收敛阶。例如在已知 $p$ 的情形，
\begin{equation}
S \approx \frac{h_2^p S(h_1) - h_1^p S(h_2)}{h_2^p - h_1^p}.
\end{equation}
扩展到多点时，可构造外推表以递推消元更高阶误差项。特点：
\begin{itemize}
  \item 对“具有明确阶数 $p$ 的渐近误差”效果好；
  \item 对不满足光滑渐近展开、或含振荡/非解析误差的序列可能不稳定；
  \item 需要对有效阶数 $p$ 有一定了解（或采用自适应估计）。
\end{itemize}

在实际应用中，若缺乏可靠的统计误差模型，一个常见的保守不确定度估计是选取参考项 $S_{\mathrm{ref}}$（例如序列最后一项或某一指定近似）并取
\begin{equation}
U=|S-S_{\mathrm{ref}}|.
\end{equation}
该 $U$ 反映的是“外推/加速值与参考项的差异”，并不必然等同于统计意义上的标准差；应结合序列差分与多种方法交叉验证。

\subsection{Shanks 变换与 Wynn $\varepsilon$ 算法}

Shanks 变换与 Wynn $\varepsilon$ 算法属于非线性序列变换，常用于加速线性收敛或对数收敛序列。其直观动机是用有理近似（Pad\'e 逼近的思想）去逼近序列的极限。

最简单的加速形式与 Aitken 的 $\Delta^2$ 过程相关：
\begin{equation}
s'_n = s_n - \frac{(\Delta s_n)^2}{\Delta^2 s_n},\quad
\Delta s_n=s_{n+1}-s_n,\quad
\Delta^2 s_n=s_{n+2}-2s_{n+1}+s_n,
\end{equation}
在 $\Delta^2 s_n$ 不接近零时可显著加速线性收敛。Wynn $\varepsilon$ 算法提供了更系统的递推表构造，实际使用时通常直接基于递推实现而非手写公式。

特点：
\begin{itemize}
  \item 对线性收敛序列经常有效，且与 Pad\'e 近似有紧密联系；
  \item 可能出现除零或数值放大（当差分很小或序列噪声较大时）；
  \item 对强振荡序列未必优于 Levin 类方法。
\end{itemize}

\subsection{Levin 变换（以 $u$-transform 为代表）}

Levin 类变换通过引入“余项估计” $\omega_n$ 来构造对极限的更好估计，常用于缓慢收敛或交错振荡的序列。其一般形式可写为
\begin{equation}
S \approx \frac{\sum_{k=0}^{n} w_k\,\frac{s_k}{\omega_k}}{\sum_{k=0}^{n} w_k\,\frac{1}{\omega_k}},
\end{equation}
其中权重 $w_k$ 与所选变体（$u,t,v$ 等）有关，$\omega_k$ 通常由差分（如 $s_{k+1}-s_k$）或启发式模型给出。

特点：
\begin{itemize}
  \item 对交错/缓慢收敛序列往往优于 Shanks/Wynn；
  \item 依赖 $\omega_k$ 的质量：若余项估计退化（例如连续项相同导致“零权重”），需要额外处理；
  \item 仍可能对噪声敏感，应结合诊断量与对比分析使用。
\end{itemize}

\section{误差传递方法（Uncertainty Propagation）}

\subsection{一阶（线性化）误差传递}

设输出为
\begin{equation}
y=f(x_1,\ldots,x_m),
\end{equation}
在输入不确定度较小、且 $f$ 在工作点附近光滑时，可对 $f$ 做一阶泰勒展开并得到标准不确定度的近似：
\begin{equation}
\sigma_y^2 \approx \sum_{i=1}^{m}\left(\frac{\partial f}{\partial x_i}\sigma_{x_i}\right)^2.
\end{equation}
若输入量存在相关性，需使用协方差矩阵形式：
\begin{equation}
\sigma_y^2 \approx \bm{g}^\mathsf{T}\,\bm{\Sigma}\,\bm{g},\qquad
\bm{g}=\nabla f,\ \bm{\Sigma}_{ij}=\mathrm{Cov}(x_i,x_j).
\end{equation}
常见特点：
\begin{itemize}
  \item 适用于“小不确定度 + 光滑函数”的局部近似；
  \item 对强非线性、非高斯、或大不确定度情形可能低估误差，需要更高阶或 Monte Carlo 方法。
\end{itemize}

\subsection{偏导数的获取：解析/符号与数值差分}

一阶误差传递依赖偏导 $\partial f/\partial x_i$。理论上可采用：
\begin{itemize}
  \item 解析偏导：手工推导或符号计算（symbolic differentiation）得到精确表达；
  \item 数值偏导：当解析难以获得时，用有限差分近似。
\end{itemize}

中心差分是常用的二阶精度近似：
\begin{equation}
\frac{\partial f}{\partial x}\bigg|_{x=x_0}\approx \frac{f(x_0+h)-f(x_0-h)}{2h}.
\end{equation}
步长 $h$ 的选择需在“截断误差（随 $h$ 变小）”与“舍入误差（随 $h$ 变小而放大）”之间折衷。对二阶中心差分，经验上可取
\begin{equation}
h \sim \varepsilon^{1/3}\,\max(1,|x_0|),
\end{equation}
其中 $\varepsilon$ 表示数值系统的相对机器精度；在多精度计算中，$\varepsilon$ 随有效位数变化，从而 $h$ 也应随精度自适应调整。

\subsection{二阶（Taylor）误差传递：均值修正与方差的二阶项}

当 $f$ 的非线性不能忽略、且输入可近似视为\textbf{独立正态}随机变量
$X_i\sim\mathcal{N}(\mu_i,\sigma_i^2)$ 时，可在工作点 $\bm{\mu}$ 附近作二阶展开：
\begin{equation}
f(\bm{X})\approx f(\bm{\mu})+\bm{g}^\mathsf{T}\Delta\bm{x}
\;+\;\frac{1}{2}\Delta\bm{x}^\mathsf{T}H\,\Delta\bm{x},
\qquad \Delta\bm{x}=\bm{X}-\bm{\mu}.
\end{equation}
其中 $\bm{g}=\nabla f$，$H$ 为 Hessian 矩阵（$H_{ij}=\partial^2 f/\partial x_i\partial x_j$）。

在该近似下，输出的\textbf{均值}与\textbf{方差}可写为
\begin{align}
\mathbb{E}[f(\bm{X})] &\approx f(\bm{\mu})+\frac{1}{2}\sum_{i=1}^{m} H_{ii}\,\sigma_i^2,\\
\mathrm{Var}[f(\bm{X})] &\approx \bm{g}^\mathsf{T}\bm{\Sigma}\bm{g}
\;+\;\frac{1}{2}\,\mathrm{tr}\!\left((H\bm{\Sigma})^2\right).
\end{align}
当 $\bm{\Sigma}=\mathrm{diag}(\sigma_1^2,\ldots,\sigma_m^2)$（独立假设）时，
\begin{equation}
\mathrm{tr}\!\left((H\bm{\Sigma})^2\right)=\sum_{i=1}^{m}\sum_{j=1}^{m} H_{ij}^2\,\sigma_i^2\sigma_j^2.
\end{equation}

特点：
\begin{itemize}
  \item 相比一阶线性化，二阶方法既能补偿非线性导致的均值偏移，也能加入 Hessian 对方差的二阶贡献；
  \item 该公式依赖“正态 + 小不确定度 + 光滑函数”等假设；若分布非高斯或函数存在定义域/不连续，建议使用 Monte Carlo。
\end{itemize}

\subsection{Monte Carlo 误差传递}

Monte Carlo 方法不显式依赖导数，而是根据输入量的概率分布进行随机抽样：
\begin{equation}
X_i^{(k)}\sim \mathcal{N}(\mu_i,\sigma_i^2),\quad
Y^{(k)} = f\!\left(X_1^{(k)},\ldots,X_m^{(k)}\right),\quad k=1,\ldots,N.
\end{equation}
随后用样本均值与样本标准差估计输出的期望与不确定度：
\begin{equation}
\bar{Y}=\frac{1}{N}\sum_{k=1}^{N}Y^{(k)},\qquad
s_Y=\sqrt{\frac{1}{N-1}\sum_{k=1}^{N}\left(Y^{(k)}-\bar{Y}\right)^2}.
\end{equation}

特点：
\begin{itemize}
  \item 可自然捕获高阶效应、非线性与截断/定义域带来的分布形变；
  \item 计算成本与样本数 $N$ 成正比；实践中可用固定随机种子复现实验；
  \item 若采样命中函数定义域外（如 $x<0$ 时 $\sqrt{x}$），对应样本会失败，需要丢弃或改用更合适的输入分布模型。
\end{itemize}

\subsection{求根的不确定度传播}

对于隐式求根问题，DataLab 支持独立输入的一阶隐函数定理传播。Monte Carlo 传播会对带不确定度的常数和数据输入进行抽样，对每个抽样问题重新求根，并报告根样本的标准差。标量二阶传播是局部近似，适用于光滑的标量实根；对于方程组、分支敏感方程或根数量可能变化的问题，优先使用 Monte Carlo。

\section{拟合方法（Fitting）}

\subsection{最小二乘与加权 $\chi^2$}

给定数据点 $\{(x_i,y_i)\}_{i=1}^{n}$ 与模型 $y_{\mathrm{model}}(x;\bm{p})$，最小二乘拟合通过最小化残差平方和：
\begin{equation}
\mathrm{RSS}(\bm{p})=\sum_{i=1}^{n} r_i(\bm{p})^2,\qquad r_i=y_{\mathrm{model}}(x_i;\bm{p})-y_i.
\end{equation}
当每个点具有标准不确定度 $\sigma_i$ 且可视为独立高斯噪声时，常用加权 $\chi^2$：
\begin{equation}
\chi^2(\bm{p})=\sum_{i=1}^{n} w_i\,r_i(\bm{p})^2,\qquad w_i=\frac{1}{\sigma_i^2}.
\end{equation}
自由度 $\nu=n-k$（$k$ 为自由参数数），缩减卡方 $\chi^2_\nu=\chi^2/\nu$ 常用于判断误差模型是否与数据一致（理想情况下 $\chi^2_\nu\approx 1$）。

\subsection{线性模型：解析解与协方差}

当模型对参数线性：
\begin{equation}
y \approx \sum_{j=1}^{k} p_j\,\varphi_j(x),
\end{equation}
记设计矩阵 $X_{ij}=\varphi_j(x_i)$，则加权最小二乘的正规方程解为
\begin{equation}
\hat{\bm{p}}=(X^\mathsf{T}WX)^{-1}X^\mathsf{T}W\bm{y},
\end{equation}
其中 $W=\mathrm{diag}(w_1,\ldots,w_n)$。在经典近似下，参数协方差为
\begin{equation}
\mathrm{Cov}(\hat{\bm{p}})\approx \sigma^2 (X^\mathsf{T}WX)^{-1},\qquad
\sigma^2\approx \frac{\chi^2}{\nu}.
\end{equation}
参数标准不确定度为 $\sigma(p_j)=\sqrt{\mathrm{Cov}_{jj}}$；非对角元反映参数相关性。

\subsection{非线性模型：局部线性化与雅可比矩阵}

对非线性模型，可在最优点附近对残差做线性化：
\begin{equation}
\bm{r}(\bm{p}+\Delta\bm{p})\approx \bm{r}(\bm{p}) + J\,\Delta\bm{p},\qquad
J_{ij}=\frac{\partial r_i}{\partial p_j}.
\end{equation}
由此得到 Gauss--Newton 近似下的协方差估计：
\begin{equation}
\mathrm{Cov}(\hat{\bm{p}})\approx \sigma^2 (J^\mathsf{T}WJ)^{-1}.
\end{equation}
该公式与“对模型输出的雅可比 $\partial y/\partial p$”等价（差一个符号），是许多非线性拟合不确定度评估的基础。

\subsection{模型比较指标：AIC / BIC / $R^2$}

为在不同模型之间权衡拟合优度与复杂度，可使用信息准则。以 $\mathrm{RSS}$（或未加权残差平方和）为噪声尺度的简化形式：
\begin{align}
\mathrm{AIC} &\approx 2k + n\ln(\mathrm{RSS}/n),\\
\mathrm{BIC} &\approx k\ln n + n\ln(\mathrm{RSS}/n).
\end{align}
$R^2$（决定系数）常用于衡量解释方差比例：
\begin{equation}
R^2 = 1-\frac{\sum_i r_i^2}{\sum_i (y_i-\bar{y})^2},
\end{equation}
但在非线性、无截距或加权情形需谨慎解释。

\section{统计平均方法（Statistics）}

\subsection{算术平均与标准误差}

算术平均：
\begin{equation}
\bar{x}=\frac{1}{n}\sum_{i=1}^{n}x_i.
\end{equation}
样本方差（无偏估计）：
\begin{equation}
s^2=\frac{1}{n-1}\sum_{i=1}^{n}(x_i-\bar{x})^2,\qquad (n>1)
\end{equation}
均值的标准误差（standard error）：
\begin{equation}
\sigma_{\bar{x}}=\frac{s}{\sqrt{n}}.
\end{equation}

\subsection{加权平均（权重 $w_i=1/\sigma_i^2$）}

当每个点具有已知标准不确定度 $\sigma_i$（独立）时，加权平均为
\begin{equation}
\bar{x}_w=\frac{\sum_i w_i x_i}{\sum_i w_i},\qquad w_i=\frac{1}{\sigma_i^2}.
\end{equation}
其均值标准不确定度为
\begin{equation}
\sigma_{\bar{x}_w}=\sqrt{\frac{1}{\sum_i w_i}}.
\end{equation}
此外常用“有效样本量”衡量权重集中程度：
\begin{equation}
n_{\mathrm{eff}}=\frac{\left(\sum_i w_i\right)^2}{\sum_i w_i^2},
\end{equation}
当权重差异很大时，$n_{\mathrm{eff}}$ 可能远小于 $n$。

若希望用同一套权重给出“数据散布”的估计，一种常见的加权样本方差无偏形式为
\begin{equation}
s_w^2=\frac{\sum_i w_i(x_i-\bar{x}_w)^2}{\sum_i w_i-\frac{\sum_i w_i^2}{\sum_i w_i}},
\end{equation}
当所有权重相等时退化为普通样本方差 $s^2$；并且该分母也可写为 $W(1-1/n_{\mathrm{eff}})$（其中 $W=\sum_i w_i$）。

\section{常见实践建议（与方法无关）}
\begin{itemize}
  \item 外推/加速结果建议与原序列趋势、差分量和多种方法交叉验证；
  \item 误差传递若存在显著相关性，应提供协方差矩阵或改用 Monte Carlo；
  \item 拟合不确定度依赖模型正确性与噪声假设；出现强相关参数时应关注协方差/相关系数；
  \item 加权统计需检查 $\sigma_i$ 的可信度，避免将系统误差当作纯统计误差处理。
\end{itemize}

\end{document}
